\begin{studiumworked}{M03. Dos rutas para la velocidad de una expansión · Intermedio}
\textbf{Enunciado.} Un gas académico de $\gamma=1.3$ y masa molar $23.3\,\mathrm{kg/kmol}$ tiene $T_0=2430$ K y $p_0=2.50$ MPa. Una sección isentrópica alcanza $M=2.075$. Halle $R,c_p,T,p,u$ y verifique energía.

\textbf{Datos.} Exponente constante: $1.3$; Masa molar: $23.3\,\mathrm{kg/kmol}$; Temperatura total: $2430\,\mathrm{K}$; Presión total: $2.500000\times10^{6}\,\mathrm{Pa}$; Mach: $2.075$.

\textbf{Incógnitas.} Propiedades específicas y estado estático.

\textbf{Hipótesis.} Gas caloríficamente perfecto de composición fija. Expansión adiabática reversible, sin trabajo ni altura.

\textbf{Modelo físico.} Energía y relaciones de remanso.
\[R=R_u/\mathcal M,\quad c_p=\gamma R/(\gamma-1)\]
\[T=T_0/[1+(\gamma-1)M^2/2]\]
\[p=p_0(T/T_0)^{\gamma/(\gamma-1)}\]
\[u=M\sqrt{\gamma RT}=\sqrt{2c_p(T_0-T)}\]
\textbf{Desarrollo.}
\begin{enumerate}
\item Se usa $R_u=8314.462618$ J/(kmol K), coherente con kg/kmol. Así $R=356.8439$ J/(kg K) y $c_p=1546.324$ J/(kg K).
\item El factor de remanso es $1.645844$. Dividir $T_0$ da $1476.446\,\mathrm{K}$; la potencia isentrópica da $288.5732\,\mathrm{kPa}$.
\item Con la temperatura estática, la ruta acústica da $1717.267\,\mathrm{m/s}$. La ruta energética usa el descenso $c_p(T_0-T)$ y produce el mismo valor. Usar $T_0$ dentro del sonido de la sección daría otra rapidez incompatible.
\end{enumerate}
\textbf{Operaciones para contrastar el cálculo.}
\begin{enumerate}

\item Magnitud 1.

\[\fitmath{y_{1}=\frac{8314.462618}{23.3}\approx 356.8439\,\mathrm{J/(kg\,K)}}\]

\item Magnitud 2.

\[\fitmath{y_{2}=2.075\cdot \sqrt{1.3\cdot 356.8438892\cdot 1476.446352}\approx 1717.267\,\mathrm{m/s}}\]

\item Magnitud 3.

\[\fitmath{y_{3}=1546.32352\cdot 1476.446352+\frac{{1717.266685}^{2}}{2}\approx 3757566\,\mathrm{J/kg}}\]

\end{enumerate}

\textbf{Resultados.} $T=1476.446$ K; $p=2.885732\times10^{5}$ Pa; $u=1717.267$ m/s.

\textbf{Interpretación.} El gas acelera y enfría, pero su temperatura total permanece fija. La presión total se mantiene sólo porque se ha supuesto isentropía.

\textbf{Comprobación.} $c_pT+u^2/2=c_pT_0=3.757566\times10^{6}$ J/kg. La entropía calculada mediante $c_p\ln(T/T_0)-R\ln(p/p_0)$ es cero dentro de redondeo.

\textbf{Límites.} No se modela química ni variación calorífica; gamma es un dato académico, no una propiedad universal de productos reales.
\end{studiumworked}
